% Generated by experiments/make_paper_numbers.py. Do not edit.
\begin{figure}[t]
\centering
\begin{tikzpicture}[x=1cm, y=1cm]
  \draw[gray!35] (0,0.0) -- (10.35,0.0);
  \node[font=\tiny, anchor=east] at (-0.12,0.0) {0\%};
  \draw[gray!35] (0,1.25) -- (10.35,1.25);
  \node[font=\tiny, anchor=east] at (-0.12,1.25) {25\%};
  \draw[gray!35] (0,2.5) -- (10.35,2.5);
  \node[font=\tiny, anchor=east] at (-0.12,2.5) {50\%};
  \draw[gray!35] (0,3.75) -- (10.35,3.75);
  \node[font=\tiny, anchor=east] at (-0.12,3.75) {75\%};
  \draw[gray!35] (0,5.0) -- (10.35,5.0);
  \node[font=\tiny, anchor=east] at (-0.12,5.0) {100\%};
  \draw[->, thick] (0,0) -- (10.5,0);
  \draw[->, thick] (0,0) -- (0,5.45);
  \draw (0.0,0) -- (0.0,-0.1) node[below, font=\tiny] {124};
  \draw (4.175,0) -- (4.175,-0.1) node[below, font=\tiny] {355};
  \draw (5.487,0) -- (5.487,-0.1) node[below, font=\tiny] {494};
  \draw (7.269,0) -- (7.269,-0.1) node[below, font=\tiny] {774};
  \draw (8.664,0) -- (8.664,-0.1) node[below, font=\tiny] {1100};
  \draw (10.0,0) -- (10.0,-0.1) node[below, font=\tiny] {1540};
  \node[font=\small] at (5,-0.72) {parameters (M), logarithmic};
  \draw[line width=1.1pt, knowInk] (0.0,1.25) -- (4.175,2.71) -- (5.487,3.405) -- (7.269,3.96) -- (8.664,4.86) -- (10.0,4.93);
  \fill[knowInk] (0.0,1.25) circle (2pt);
  \fill[knowInk] (4.175,2.71) circle (2pt);
  \fill[knowInk] (5.487,3.405) circle (2pt);
  \fill[knowInk] (7.269,3.96) circle (2pt);
  \fill[knowInk] (8.664,4.86) circle (2pt);
  \fill[knowInk] (10.0,4.93) circle (2pt);
  \node[font=\scriptsize, anchor=west, knowInk] at (10.6,4.93) {holds the fact};
  \draw[line width=1.1pt, adoptInk, dashed] (0.0,2.5) -- (4.175,2.18) -- (5.487,4.59) -- (7.269,3.77) -- (8.664,1.93) -- (10.0,4.365);
  \draw[adoptInk, fill=white, line width=1pt] (0.0,2.5) circle (2pt);
  \draw[adoptInk, fill=white, line width=1pt] (4.175,2.18) circle (2pt);
  \draw[adoptInk, fill=white, line width=1pt] (5.487,4.59) circle (2pt);
  \draw[adoptInk, fill=white, line width=1pt] (7.269,3.77) circle (2pt);
  \draw[adoptInk, fill=white, line width=1pt] (8.664,1.93) circle (2pt);
  \draw[adoptInk, fill=white, line width=1pt] (10.0,4.365) circle (2pt);
  \node[font=\scriptsize, anchor=west, adoptInk] at (10.6,4.365) {adopts the document};
  \draw[gray!70, dash pattern=on 1.5pt off 1.5pt] (8.664,1.93) -- (8.664,4.86);
  \node[font=\tiny, gray!70, anchor=north] at (8.664,1.79) {TinyLlama-1.1B-Chat};
\end{tikzpicture}
\caption{Holding a fact scales with parameter count; deferring to a document that contradicts it does not. Filled markers are how often the model ranks the true answer first with no document present; open markers are how often it ranks the document's answer first, over the conflicts only, so the second curve is conditioned on the first. The first is monotone in size with no inversion anywhere; the second is not ordered by size at all. The dashed segment marks TinyLlama-1.1B, which holds the fact second most reliably of the six and defers least. Coordinates are generated from the same measurement as Table~\ref{tab:results}.}
\label{fig:scale}
\end{figure}
